% Standalone research note for the hybrid-local-solver project.
%
% This document is a curated reconstruction of the project conversation on
% active-volume flattening, spider obstructions, safe RPPR support gates, and
% acceleration over expanding subspaces.  It is not a verbatim transcript.
% Source results, new proved statements, conditional implications, refuted
% claims, and open lemmas are labeled separately.

\documentclass[11pt]{article}

% Common class-neutral preamble for standalone notes under manuscript/notes/.
% Note entry points all live exactly two directories below manuscript/.

\usepackage[margin=1in]{geometry}
\usepackage{amsmath,amssymb,amsthm,mathtools,bm}
\usepackage{algorithm}
\usepackage{algorithmic}
\usepackage{booktabs,tabularx,multirow,array,longtable}
\usepackage{enumitem}
\usepackage{microtype}
\usepackage[numbers,sort&compress]{natbib}
\usepackage{xcolor}
\usepackage[colorlinks=true,citecolor=blue,linkcolor=black,urlcolor=blue]{hyperref}
\usepackage[capitalize,nameinlink,noabbrev]{cleveref}

\newtheorem{theorem}{Theorem}[section]
\newtheorem{lemma}[theorem]{Lemma}
\newtheorem{proposition}[theorem]{Proposition}
\newtheorem{corollary}[theorem]{Corollary}
\newtheorem{conjecture}[theorem]{Conjecture}
\theoremstyle{definition}
\newtheorem{definition}[theorem]{Definition}
\newtheorem{assumption}[theorem]{Assumption}
\newtheorem{problem}[theorem]{Problem}
\newtheorem{observation}[theorem]{Observation}
\theoremstyle{remark}
\newtheorem{remark}[theorem]{Remark}
\theoremstyle{plain}

% Canonical mathematical notation for the active manuscript.
%
% This file preserves the recurring notation style used in the author's
% NeurIPS 2024 and 2025 papers while excluding unrelated tensor and
% machine-learning boilerplate from those archived command collections.
%
% Prefix convention:
%   \v...  bold vectors
%   \m...  bold matrices
%   \g...  calligraphic sets, graphs, and families
%   \s...  blackboard-bold spaces and number systems

\newcommand{\method}{\textsc{HybridLocalSolver}}

% Font and scalar shorthands retained from the archived manuscripts.
\newcommand{\mc}{\mathcal}
\newcommand{\R}{\mathbb{R}}
\newcommand{\E}{\mathbb{E}}
\newcommand{\G}{\mathbb{G}}
\newcommand{\N}{\mathcal{N}}
% Accuracy namespaces are semantic: do not add a bare \eps alias.
% Degree-normalized residual tolerance of classical APPR is kept distinct from
% the objective-gap and proximal fixed-point tolerances of the RPPR sections.
\newcommand{\epsappr}{\varepsilon_{\mathrm{appr}}}
\newcommand{\epsobj}{\varepsilon_{\mathrm{obj}}}
\newcommand{\epspg}{\varepsilon_{\mathrm{pg}}}
\newcommand{\epsppr}{\varepsilon_{\mathrm{ppr}}}
\newcommand{\epsin}{\varepsilon_{\mathrm{in}}}
\newcommand{\epskkt}{\varepsilon_{\mathrm{kkt}}}
\newcommand{\epsburn}{\varepsilon_{\mathrm{burn}}}
\newcommand{\epssol}{\varepsilon_{\mathrm{sol}}}
\newcommand{\epspprinst}[1]{\varepsilon_{\mathrm{ppr},#1}}
\newcommand{\trans}{\mathsf{T}}
\newcommand{\grad}{\nabla}

% Delimiters and generic helpers.
% Norms and inner products use fixed-size delimiters by default.
% Use an optional size (e.g. \norm[\big]{...}) or * for automatic sizing.
\DeclarePairedDelimiter{\norm}{\lVert}{\rVert}
\newcommand{\abs}[1]{\left\lvert #1 \right\rvert}
\DeclarePairedDelimiterX{\inner}[2]{\langle}{\rangle}{#1, #2}
\newcommand{\ip}[2]{\inner{#1}{#2}}
\newcommand{\card}[1]{\left\lvert #1 \right\rvert}
\newcommand{\set}[1]{\left\{#1\right\}}
\newcommand{\ceil}[1]{\left\lceil #1 \right\rceil}
\newcommand{\floor}[1]{\left\lfloor #1 \right\rfloor}
\newcommand{\comp}[1]{\overline{#1}}
\newcommand{\conj}[1]{\overline{#1}}
\newcommand{\mean}[1]{\overline{#1}}
\newcommand{\bigO}{\mathcal{O}}
\newcommand{\softO}{\widetilde{\mathcal{O}}}
\newcommand{\softOmega}{\widetilde{\Omega}}
\newcommand{\pos}[1]{\left[#1\right]_{+}}

% Bold lowercase vectors.
\newcommand{\va}{\bm{a}}
\newcommand{\vb}{\bm{b}}
\newcommand{\vc}{\bm{c}}
\newcommand{\vd}{\bm{d}}
\newcommand{\ve}{\bm{e}}
\newcommand{\vf}{\bm{f}}
\newcommand{\vg}{\bm{g}}
\newcommand{\vh}{\bm{h}}
\newcommand{\vi}{\bm{i}}
\newcommand{\vj}{\bm{j}}
\newcommand{\vk}{\bm{k}}
\newcommand{\vl}{\bm{l}}
\newcommand{\vm}{\bm{m}}
\newcommand{\vn}{\bm{n}}
\newcommand{\vo}{\bm{o}}
\newcommand{\vp}{\bm{p}}
\newcommand{\vq}{\bm{q}}
\newcommand{\vr}{\bm{r}}
\newcommand{\vs}{\bm{s}}
\newcommand{\vt}{\bm{t}}
\newcommand{\vu}{\bm{u}}
\newcommand{\vv}{\bm{v}}
\newcommand{\vw}{\bm{w}}
\newcommand{\vx}{\bm{x}}
\newcommand{\vy}{\bm{y}}
\newcommand{\vz}{\bm{z}}

% Bold Greek vectors and special vectors.
\newcommand{\valpha}{\bm{\alpha}}
\newcommand{\vbeta}{\bm{\beta}}
\newcommand{\vdelta}{\bm{\delta}}
\newcommand{\vepsilon}{\bm{\epsilon}}
\newcommand{\veta}{\bm{\eta}}
\newcommand{\vgamma}{\bm{\gamma}}
\newcommand{\vlambda}{\bm{\lambda}}
\newcommand{\vmu}{\bm{\mu}}
\newcommand{\vomega}{\bm{\omega}}
\newcommand{\vphi}{\bm{\phi}}
\newcommand{\vpi}{\bm{\pi}}
\newcommand{\vpsi}{\bm{\psi}}
\newcommand{\vrho}{\bm{\rho}}
\newcommand{\vsigma}{\bm{\sigma}}
\newcommand{\vtheta}{\bm{\theta}}
\newcommand{\vxi}{\bm{\xi}}
\newcommand{\vell}{\bm{\ell}}
\newcommand{\vDelta}{\bm{\Delta}}
\newcommand{\vGamma}{\bm{\Gamma}}
\newcommand{\vLambda}{\bm{\Lambda}}
\newcommand{\vPhi}{\bm{\Phi}}
\newcommand{\vSigma}{\bm{\Sigma}}
\newcommand{\vzero}{\bm{0}}
\newcommand{\vone}{\bm{1}}
\newcommand{\one}{\mathbf{1}}
\newcommand{\zero}{\vzero}
\newcommand{\eunit}[1]{\bm{e}_{#1}}
\newcommand{\vDeltax}{\bm{\Delta}\bm{x}}

% Bold uppercase matrices.
\newcommand{\mA}{\bm{A}}
\newcommand{\mB}{\bm{B}}
\newcommand{\mC}{\bm{C}}
\newcommand{\mD}{\bm{D}}
\newcommand{\mE}{\bm{E}}
\newcommand{\mF}{\bm{F}}
\newcommand{\mG}{\bm{G}}
\newcommand{\mH}{\bm{H}}
\newcommand{\mI}{\bm{I}}
\newcommand{\mJ}{\bm{J}}
\newcommand{\mK}{\bm{K}}
\newcommand{\mL}{\bm{L}}
\newcommand{\mM}{\bm{M}}
\newcommand{\mN}{\bm{N}}
\newcommand{\mO}{\bm{O}}
\newcommand{\mP}{\bm{P}}
\newcommand{\mQ}{\bm{Q}}
\newcommand{\mR}{\bm{R}}
\newcommand{\mS}{\bm{S}}
\newcommand{\mT}{\bm{T}}
\newcommand{\mU}{\bm{U}}
\newcommand{\mV}{\bm{V}}
\newcommand{\mW}{\bm{W}}
\newcommand{\mX}{\bm{X}}
\newcommand{\mY}{\bm{Y}}
\newcommand{\mZ}{\bm{Z}}

% Bold Greek matrices retained from the graph papers.
\newcommand{\mBeta}{\bm{\beta}}
\newcommand{\mDelta}{\bm{\Delta}}
\newcommand{\mGamma}{\bm{\Gamma}}
\newcommand{\mLambda}{\bm{\Lambda}}
\newcommand{\mOmega}{\bm{\Omega}}
\newcommand{\mPhi}{\bm{\Phi}}
\newcommand{\mPi}{\bm{\Pi}}
\newcommand{\mSigma}{\bm{\Sigma}}
\newcommand{\mTheta}{\bm{\Theta}}

% Calligraphic symbols.
\newcommand{\gA}{\mathcal{A}}
\newcommand{\gB}{\mathcal{B}}
\newcommand{\gC}{\mathcal{C}}
\newcommand{\gD}{\mathcal{D}}
\newcommand{\gE}{\mathcal{E}}
\newcommand{\gF}{\mathcal{F}}
\newcommand{\gG}{\mathcal{G}}
\newcommand{\gH}{\mathcal{H}}
\newcommand{\gI}{\mathcal{I}}
\newcommand{\gJ}{\mathcal{J}}
\newcommand{\gK}{\mathcal{K}}
\newcommand{\gL}{\mathcal{L}}
\newcommand{\gM}{\mathcal{M}}
\newcommand{\gN}{\mathcal{N}}
\newcommand{\gO}{\mathcal{O}}
\newcommand{\gP}{\mathcal{P}}
\newcommand{\gQ}{\mathcal{Q}}
\newcommand{\gR}{\mathcal{R}}
\newcommand{\gS}{\mathcal{S}}
\newcommand{\gT}{\mathcal{T}}
\newcommand{\gU}{\mathcal{U}}
\newcommand{\gV}{\mathcal{V}}
\newcommand{\gW}{\mathcal{W}}
\newcommand{\gX}{\mathcal{X}}
\newcommand{\gY}{\mathcal{Y}}
\newcommand{\gZ}{\mathcal{Z}}

% Blackboard-bold symbols.
\newcommand{\sA}{\mathbb{A}}
\newcommand{\sB}{\mathbb{B}}
\newcommand{\sC}{\mathbb{C}}
\newcommand{\sD}{\mathbb{D}}
\newcommand{\sE}{\mathbb{E}}
\newcommand{\sF}{\mathbb{F}}
\newcommand{\sG}{\mathbb{G}}
\newcommand{\sH}{\mathbb{H}}
\newcommand{\sI}{\mathbb{I}}
\newcommand{\sJ}{\mathbb{J}}
\newcommand{\sK}{\mathbb{K}}
\newcommand{\sL}{\mathbb{L}}
\newcommand{\sM}{\mathbb{M}}
\newcommand{\sN}{\mathbb{N}}
\newcommand{\sO}{\mathbb{O}}
\newcommand{\sP}{\mathbb{P}}
\newcommand{\sQ}{\mathbb{Q}}
\newcommand{\sR}{\mathbb{R}}
\newcommand{\sS}{\mathbb{S}}
\newcommand{\sT}{\mathbb{T}}
\newcommand{\sU}{\mathbb{U}}
\newcommand{\sV}{\mathbb{V}}
\newcommand{\sW}{\mathbb{W}}
\newcommand{\sX}{\mathbb{X}}
\newcommand{\sY}{\mathbb{Y}}
\newcommand{\sZ}{\mathbb{Z}}

% Frequently used operators.
\DeclareMathOperator{\Cov}{Cov}
\DeclareMathOperator{\Var}{Var}
\DeclareMathOperator{\diag}{diag}
\DeclareMathOperator{\nei}{Nei}
\DeclareMathOperator{\nnz}{nnz}
\DeclareMathOperator{\pr}{pr}
\DeclareMathOperator{\prox}{prox}
\DeclareMathOperator{\push}{push}
\DeclareMathOperator{\res}{res}
\DeclareMathOperator{\rel}{Rel}
\DeclareMathOperator{\relu}{ReLU}
\DeclareMathOperator{\sign}{sign}
\DeclareMathOperator{\supp}{supp}
\DeclareMathOperator{\tr}{tr}
\DeclareMathOperator{\Tr}{Tr}
\DeclareMathOperator{\vol}{vol}
\DeclareMathOperator{\work}{work}
\DeclareMathOperator{\Work}{Work}
\DeclareMathOperator*{\argmax}{arg\,max}
\DeclareMathOperator*{\argmin}{arg\,min}

% Project-wide names and claim-status labels used by standalone research notes.
% Mathematical notation belongs in math_commands.tex.

% Algorithms and solver families.
\newcommand{\AESP}{\textsc{AESP}}
\newcommand{\APPR}{\textsc{APPR}}
\newcommand{\ASPR}{\textsc{ASPR}}
\newcommand{\FISTA}{\textsc{FISTA}}
\newcommand{\ISTA}{\textsc{ISTA}}
\newcommand{\LOCAPPR}{\textsc{LocAPPR}}
\newcommand{\LOCGD}{\textsc{LocGD}}
\newcommand{\LOCSOR}{\textsc{LocSOR}}
\newcommand{\LOCCH}{\textsc{LocCH}}
\newcommand{\LOCHB}{\textsc{LocHB}}
\newcommand{\LocProxGS}{\textsc{LocProxGS}}
\newcommand{\Hybrid}{\textsc{AESP--LocSOR}}

% Claim classifications required by the repository research-note protocol.
\newcommand{\statussource}{\textbf{Source result}}
\newcommand{\statusproved}{\textbf{Proved here}}
\newcommand{\statusconditional}{\textbf{Conditional}}
\newcommand{\statusconjecture}{\textbf{Open / conjectural}}
\newcommand{\statusopen}{\textbf{Open}}
\newcommand{\statusempirical}{\textbf{Empirical}}
\newcommand{\statusobservation}{\textbf{Empirical observation}}
\newcommand{\statusrefuted}{\textbf{Refuted route}}

% Compatibility names used by the older complete-note presentation layer.
\newcommand{\Status}[2]{\textbf{#1:} #2}
\newcommand{\SourceStatus}{\textnormal{\textsc{Source result}}}
\newcommand{\ProvedStatus}{\textnormal{\textsc{Proved here}}}
\newcommand{\ConditionalStatus}{\textnormal{\textsc{Conditional}}}
\newcommand{\ComputedStatus}{\textnormal{\textsc{Computed}}}
\newcommand{\RefutedStatus}{\textnormal{\textsc{Refuted route}}}
\newcommand{\OpenStatus}{\textnormal{\textsc{Open}}}



\title{Volume-Gated Acceleration for Local PageRank:\\
Spectral Structure, Safe Support Certification, and the Expanding-Subspace Lemma}
\author{Baojian Zhou\\
\small Working research note for \texttt{hybrid-local-solver}}
\date{August 2026}

\begin{document}
\maketitle

\begin{abstract}
Accelerated local PageRank methods can reduce the number of iteration-scale
steps while enlarging the degree volume scanned by each step.  This note asks
whether one can flatten the observed bell-shaped active-volume curve: permit
$\bigO(\alpha^{-1/2}\log(1/\epsppr))$ iterations, but keep every active set
at volume $\bigO(1/\epsppr)$.  We first separate conditioning from locality.
Every principal PageRank system has condition number at most $1/\alpha$, and
the exact condition number of a depth-$L$ spider prefix is
$\Theta(1/(\alpha+L^{-2}))$.  Thus the spider obstruction comes from support
growth and edge revisits, not a fictitious local condition number
$\Theta(\alpha^{-2})$.

We then prove the support-certification part of a volume-gated method.  Exact
PPR has a degree-volume-$1/\tau$ superlevel core whose Dirichlet restriction is
$\tau$-accurate in degree-normalized infinity norm.  More importantly, the
$\ell_1$-regularized PageRank solution has support volume at most $1/\rho$ and
is $\rho$-accurate for unregularized PPR.  An arbitrary signed accelerated
candidate on a certified subspace can be converted into a componentwise lower
envelope.  Negative KKT violations at this envelope safely reveal only true
support vertices, while a one-sided global KKT test certifies solution error.
Taking $\rho=\tau=\epsppr/2$ therefore proves a uniform peak-volume cap
$2/\epsppr$ and final PPR error at most $\epsppr$.

The remaining acceleration claim is not closed.  We give two corrections to
the initial proof route.  First, an endpoint-source path forces
$\Omega(\alpha^{-1/2}\log(\sqrt\alpha/\rho))$ one-vertex support expansions,
so restarting after every expansion cannot yield only
$\bigO(\log(1/\epsppr))$ restarts uniformly in $\alpha$.  Second, the
Schur-complement ``shock energy'' lower bound has the wrong direction for a
stability proof.  A valid telescoping argument nevertheless bounds all
continuous re-solving work by
$\bigO(\alpha^{-1/2}\log(1/\epsppr))$ accelerated sweeps plus one discrete
term per expansion.  Closing the desired
$\softO(1/(\sqrt\alpha\epsppr))$ work theorem now reduces to a precise
one-sided, projected, expanding-subspace acceleration lemma.  The universal
work bound is explicitly left open.
\end{abstract}

\tableofcontents

\section{Scope, status, and source map}
\label{sec:scope}

This note reconstructs the complete volume-flattening discussion that grew
out of the AESP--local-refinement project.  It complements, but does not
replace, the broader standalone AESP--\LOCSOR{} note.  In particular, it does
not promote a graph-uniform end-to-end complexity theorem into the active
manuscript.

\begin{center}
\begin{tabularx}{\textwidth}{@{}p{0.18\textwidth}X@{}}
\toprule
Status & Statement \\
\midrule
\statussource & The RPPR objective, support-volume bound, classical
\FISTA{} work model, confinement analysis, and bad-star phenomenon are from
\citet{fountoulakis2026}.  The locally evolving-set and nonhomogeneous
\LOCCH{} recurrences are from \citet{zhou2024}; AESP is from
\citet{huang2025}. \\
\statusproved & Principal restricted PageRank systems have condition number
at most $1/\alpha$; a depth-$L$ spider prefix has condition number
$\Theta(1/(\alpha+L^{-2}))$. \\
\statusproved & A $1/\tau$-volume exact PPR core exists; fixed RPPR
regularization gives a $1/\rho$-volume safe support; signed candidates admit
a safe lower-envelope correction; and one-sided KKT tests give safe admission
and an infinity-error certificate. \\
\statusproved & The exact stagewise gate terminates and, with
$\rho=\tau=\epsppr/2$, returns an $\epsppr$-accurate PPR approximation while
every working set has volume at most $2/\epsppr$. \\
\statusproved & A path forces
$\Omega(\alpha^{-1/2}\log(\sqrt\alpha/\rho))$ expansions; the exact expansion
gain is a Schur-complement quadratic; and continuous restricted re-solving
costs telescope to an accelerated term plus one discrete term per expansion.
\\
\statusconditional & Geometric volume epochs would give
$\bigO(1/(\rho\sqrt\alpha))$ work if subspace discovery and continuation can
be charged without repeatedly scanning the whole current support. \\
\statusrefuted & The spider does not have local condition number
$\Theta(\alpha^{-2})$; splitting large sweeps into batches does not reduce
work; support changes cannot in general be grouped into only
$\bigO(\log(1/\epsppr))$ ordinary restarts; and a lower bound on expansion
gain cannot serve as the upper perturbation bound needed for stability. \\
\statusopen & A one-sided projected acceleration lemma for safely expanding
subspaces, or a weaker newly-admitted-volume charging lemma, is still needed
for a graph-uniform $\softO(1/(\rho\sqrt\alpha))$ work theorem. \\
\bottomrule
\end{tabularx}
\end{center}

\paragraph{Source pointers.}
The RPPR formulation and KKT conditions appear on PDF page 3, Section 3, of
\citet{fountoulakis2026}; its FISTA update and degree-weighted work are on PDF
page 4, Section 3.1; its conditional confinement theorems are on PDF pages
5--6, Sections 4.2--4.3; its star lower bound is in Appendix D, PDF pages
18--23; and its experimental fixed-point residual is on PDF page 24.  The
locally evolving-set work model and localized Chebyshev recurrence appear in
Sections 3--4 and Appendix C of \citet{zhou2024}.  The AESP active-volume
model, inexact proximal maps, and PPR specialization appear on PDF pages 4--7
and 26--29 of \citet{huang2025}.  These source locations are also recorded in
the repository literature notes.

\paragraph{Note-scoped conventions.}
The repository-wide residual convention is unresolved.  In this note:
$\rho$ is the RPPR regularization level; $\tau$ is a one-sided KKT/error
tolerance; and $\epsppr$ is only the final PPR error target
\begin{equation}
  \norm{D^{-1/2}(\widehat x-x^0)}_\infty\leq\epsppr.
  \label{eq:note-error}
\end{equation}
One full restricted first-order step on $U$ is charged $\vol(U)$, corresponding
to scanning every incident edge of every active coordinate once.  No
conversion to the repository's APPR threshold, objective-gap tolerance, or
eventual implementation stopping rule is asserted.

\section{PageRank and RPPR as a Stieltjes quadratic}
\label{sec:model}

% Canonical source-aligned PageRank/RPPR reference formulation.
%
% This fragment is intentionally heading-free at the section level so that the
% active paper and every standalone note can input it. It is a manuscript
% reference convention, not an implementation-wide residual decision.
% Primary source: Fountoulakis and Martinez-Rubio, arXiv:2602.21138v2,
% PDF p. 3, Section 3 and equation (RPPR).

\subsection{Common source-aligned PageRank and RPPR model}
\label{subsec:shared-source-aligned-problem}

All documents in this manuscript workspace use the following reference
language. It follows the plain italic notation of the comparison paper:
\(x,s,A,D,Q,I\) denote column vectors or matrices as determined by their
displayed domains. This is the scoped exception to the project's usual
bold-vector and bold-matrix typography. A note may impose a stronger
assumption after this block, but it must not redefine these objects.

Let \(G=(V,E)\) be a finite, simple, undirected, unweighted graph without
isolated vertices, let \(V=[n]:=\{1,\ldots,n\}\), and let
\(A\in\{0,1\}^{n\times n}\) be its symmetric adjacency matrix. For
\(i\in V\), write
\[
    \mathcal N(i):=\{j\in V:\{i,j\}\in E\},
    \qquad
    d_i:=|\mathcal N(i)|,
    \qquad
    D:=\diag(d_1,\ldots,d_n).
\]
For \(S\subseteq V\), define
\[
    \mathcal N(S):=\bigcup_{i\in S}\mathcal N(i),
    \qquad
    \partial S:=\mathcal N(S)\setminus S,
    \qquad
    \operatorname{Ext}(S):=V\setminus(S\cup\partial S),
\]
\begin{equation}
    \vol(S):=\sum_{i\in S}d_i.
    \label{eq:shared-volume}
\end{equation}
The support of a vector is
\(\supp(x):=\{i\in[n]:x_i\neq0\}\). Matrix inequalities use the Loewner
order, and all unqualified vector inequalities are coordinatewise.

The seed is a column distribution
\begin{equation}
    s\in\R_+^n,
    \qquad
    \inner{\one}{s}=1.
    \label{eq:shared-seed}
\end{equation}
The single-seed specialization is always written \(s=e_v\), where \(v\in V\)
is the seed vertex. Thus \(s\) is never used as a vertex label. For
\(\alpha\in(0,1]\), define
\begin{equation}
\begin{aligned}
    \mathcal L&:=I-D^{-1/2}AD^{-1/2},\\
    Q&:=\alpha I+\frac{1-\alpha}{2}\mathcal L
      =\frac{1+\alpha}{2}I
       -\frac{1-\alpha}{2}D^{-1/2}AD^{-1/2},\\
    b&:=\alpha D^{-1/2}s.
\end{aligned}
    \label{eq:shared-pagerank-matrices}
\end{equation}
Since \(0\preceq\mathcal L\preceq2I\),
\(\alpha I\preceq Q\preceq I\). The shared smooth PageRank quadratic is
\begin{equation}
    f(x):=\frac12\inner{x}{Qx}-\inner{b}{x},
    \qquad
    \nabla f(x)=Qx-b.
    \label{eq:shared-pagerank-objective}
\end{equation}
Its unique unregularized minimizer and associated lazy personalized PageRank
vector are
\begin{equation}
    x^0:=Q^{-1}b,
    \qquad
    \pi:=D^{1/2}x^0.
    \label{eq:shared-ppr-solution}
\end{equation}

For a regularization parameter \(\rho>0\), reserve \(g_\rho\) for the
nonsmooth weighted \(\ell_1\) term and define
\begin{equation}
    g_\rho(x):=\alpha\rho\norm{D^{1/2}x}_1,
    \qquad
    F_\rho(x):=f(x)+g_\rho(x),
    \qquad
    x^\star(\rho):=\argmin_{x\in\R^n}F_\rho(x).
    \label{eq:shared-rppr-objective}
\end{equation}
When \(\rho\) is fixed, \(x^\star\) abbreviates \(x^\star(\rho)\), never the
unregularized point \(x^0\). Write
\[
    S^\star(\rho):=\supp(x^\star(\rho)),
    \qquad
    I^\star(\rho):=[n]\setminus S^\star(\rho).
\]
The source records \(x^\star(\rho)\geq0\),
\(\vol(S^\star(\rho))\leq1/\rho\), and the coordinatewise conditions
\begin{equation}
    \nabla_i f(x^\star(\rho))
    \in
    \begin{cases}
        \{-\alpha\rho\sqrt{d_i}\}, & x_i^\star(\rho)>0,\\
        [-\alpha\rho\sqrt{d_i},0], & x_i^\star(\rho)=0.
    \end{cases}
    \label{eq:shared-rppr-kkt}
\end{equation}

The common computational unit is one repeated adjacency-list scan: scanning
coordinate \(i\) costs \(d_i\), and scanning a set \(S\) costs \(\vol(S)\).
Iteration count and wall-clock time do not replace this degree-weighted work
measure.

Finally, accuracy symbols remain distinct. The APPR threshold is
\(\epsappr\); \(\epsobj\) denotes an objective-gap target;
\(\epspg\) denotes the source experiment's proximal fixed-point tolerance; and
\(\epsppr\) may denote a document-scoped degree-normalized PPR error target.
No equality or conversion among these quantities is assumed. In particular,
this shared model does not resolve the repository-wide residual decision.


\paragraph{Note specialization.}
This note additionally assumes that \(G\) is connected and
\(0<\alpha<1\). The shared matrix \(Q\) is then a symmetric nonsingular
Stieltjes matrix:
\begin{equation}
  \alpha I\preceq Q\preceq I,
  \qquad Q_{ij}\leq0\ (i\neq j),
  \qquad Q^{-1}\geq0,
  \qquad QD^{1/2}\one=\alpha D^{1/2}\one.
  \label{eq:Q-properties}
\end{equation}
The shared unregularized point \(x^0\) is nonnegative. The last identity in
\cref{eq:Q-properties} gives
\begin{equation}
  \one^\top\pi=1.
  \label{eq:ppr-mass}
\end{equation}
It is useful to expose the one-sided obstacle structure of the shared RPPR
objective in \cref{eq:shared-rppr-objective}.

\begin{lemma}[Equivalent nonnegative quadratic program]
\label{lem:nonnegative-equivalence}
The unique minimizer \(x^\star(\rho)\), abbreviated here by
\(x_\rho^\star\), is nonnegative and is
equivalently the unique minimizer of
\begin{equation}
  \min_{x\geq0}G_\rho(x),
  \qquad
  G_\rho(x):=\frac12x^\top Qx-c_\rho^\top x,
  \qquad
  c_\rho:=b-\alpha\rho D^{1/2}\one.
  \label{eq:obstacle-QP}
\end{equation}
It satisfies the complementarity system
\begin{equation}
 x_\rho^\star\geq0,
 \qquad
 \kappa_\rho(x_\rho^\star):=Qx_\rho^\star-c_\rho\geq0,
 \qquad
 x_{\rho,i}^\star \kappa_{\rho,i}(x_\rho^\star)=0.
 \label{eq:KKT}
\end{equation}
\end{lemma}

\begin{proof}
Because the off-diagonal entries of $Q$ are nonpositive,
$\abs{x}^\top Q\abs{x}\leq x^\top Qx$.  Moreover
$b^\top\abs{x}\geq b^\top x$ and the weighted $\ell_1$ term is unchanged.
Thus $F_\rho(\abs{x})\leq F_\rho(x)$.  Strong convexity makes the minimizer
unique, hence nonnegative.  On the nonnegative orthant the weighted
$\ell_1$ term is linear, giving \cref{eq:obstacle-QP};
\cref{eq:KKT} is its KKT system.
\end{proof}

For a vector $v$, write $(v_-)_i:=\max\{-v_i,0\}$.  All componentwise
inequalities below use the ordinary coordinate order.

\section{The spider obstruction is geometric, not spectral}
\label{sec:spider}

The first question in the discussion was whether the local system on a long
spider might have condition number $\Theta(\alpha^{-2})$, thereby erasing the
usual square-root acceleration.  Principal-submatrix interlacing rules this
out immediately.

\begin{lemma}[Uniform restricted conditioning]
\label{lem:principal-conditioning}
For every $U\subseteq V$,
\begin{equation}
  \alpha I_U\preceq Q_{UU}\preceq I_U,
  \qquad
  \kappa_2(Q_{UU})\leq\frac1\alpha.
  \label{eq:principal-kappa}
\end{equation}
\end{lemma}

\begin{proof}
$Q_{UU}$ is a principal submatrix of $Q$.  The claim follows from
\cref{eq:Q-properties}, equivalently from Cauchy interlacing.
\end{proof}

The exact spider calculation identifies the scale more sharply.  Let
$G_{B,N}$ have a center $o$ and $B$ disjoint arms of length $N$.  Let $U_L$
contain $o$ and the first $L$ vertices of each arm, where $1\leq L<N$; degrees
are always those of the full spider, so every truncated depth-$L$ vertex has
degree two.

\begin{proposition}[Exact spider-prefix spectrum]
\label{prop:spider-spectrum}
Let $S=D^{-1/2}AD^{-1/2}$.  Then
\begin{equation}
 \lambda_{\max}(S_{U_LU_L})
 =\cos\!\left(\frac{\pi}{2(L+1)}\right).
 \label{eq:spider-lambda}
\end{equation}
Consequently, with the spider-scoped ratio
$\beta_{\mathrm{sp}}:=(1-\alpha)/(1+\alpha)$,
\begin{align}
 \kappa_2(Q_{U_LU_L})
 &=\frac{1+\beta_{\mathrm{sp}}\cos(\pi/(2(L+1)))}
         {1-\beta_{\mathrm{sp}}\cos(\pi/(2(L+1)))}
 \label{eq:spider-kappa-exact}\\
 &=\Theta\!\left(\frac1{\alpha+L^{-2}}\right)
 =\Theta\!\left(\min\left\{L^2,\frac1\alpha\right\}\right).
 \label{eq:spider-kappa-asymptotic}
\end{align}
\end{proposition}

\begin{proof}
The Perron eigenvector is invariant under permutations of the arms.  In the
orthonormal radial basis consisting of the center and the uniform vector on
each depth, $S_{U_LU_L}$ becomes a tridiagonal matrix whose first off-diagonal
entry is $1/\sqrt2$ and whose remaining off-diagonal entries are $1/2$.
Its radial eigenvalues are
\[
  \cos\!\left(\frac{(2k-1)\pi}{2(L+1)}\right),
  \qquad 1\leq k\leq L+1.
\]
The nonradial arm-difference modes have strictly smaller spectral radius, so
\cref{eq:spider-lambda} follows.  The spider is bipartite, hence the extreme
restricted normalized-adjacency eigenvalues are the positive and negative of
\cref{eq:spider-lambda}. Substitution into
\cref{eq:shared-pagerank-matrices} gives
\cref{eq:spider-kappa-exact}.  Finally,
$1-\beta_{\mathrm{sp}}\cos(\pi/(2(L+1)))
=\Theta(\alpha+L^{-2})$, whereas the numerator is
bounded above and below by constants.
\end{proof}

Thus $L=\Theta(\alpha^{-1/2})$ already saturates the global condition number
$\Theta(1/\alpha)$.  Optimal polynomial or SOR-type acceleration still needs
only $\Theta(1/\sqrt\alpha)$ sweep-scale iterations.  The difficulty is the
cost of those sweeps.

\subsection{A space--time accounting model}

Write
\begin{equation}
  B:=\Theta(1/\epsppr),
  \qquad L:=\Theta(1/\sqrt\alpha).
  \label{eq:B-L}
\end{equation}
Suppose a FIFO wavefront scans all $B$ arm prefixes through depth $j$ at its
$j$th sweep.  This is a structural accounting model, not a graph-uniform
lower bound for every SOR implementation.  It gives
\begin{equation}
 \vol(S_j)=\Theta(Bj),
 \qquad
 \sum_{j=1}^{L}\vol(S_j)
 =\Theta(BL^2)
 =\Theta\!\left(\frac1{\alpha\epsppr}\right).
 \label{eq:triangular-work}
\end{equation}
Equivalently, a late sweep costs $BL$ and there are $L$ accelerated-scale
sweeps.  Merely splitting a size-$Bj$ sweep into $j$ batches of size $B$ does
not change the sum.

\begin{center}
\begin{tabular}{@{}lccc@{}}
\toprule
Dynamics & Peak active volume & Iteration/batch count & Total work \\
\midrule
Growing FIFO prefix & $BL$ & $L$ sweeps & $BL^2$ \\
Same work split into batches & $B$ & $L^2$ batches & $BL^2$ \\
Desired volume-gated method & $B$ & $L\log(1/\epsppr)$ & $BL\log(1/\epsppr)$ \\
\bottomrule
\end{tabular}
\end{center}

The desired method must therefore freeze old interior coordinates or charge
their revisits amortizedly.  A change in iteration labels is insufficient.

\section{The conditional geometric-epoch tradeoff}
\label{sec:epoch-blueprint}

The cleanest algebraic expression of volume flattening uses geometrically
decreasing tolerances
\begin{equation}
  \tau_j=2^{-j}\tau_0,
  \qquad J=\Theta(\log(\tau_0/\tau_J)).
  \label{eq:geometric-tolerances}
\end{equation}

\begin{proposition}[Conditional geometric-epoch work]
\label{prop:conditional-epochs}
Suppose that at epoch $j$ an algorithm has a certified region $U_j$ with
\begin{equation}
  \vol(U_j)\leq\frac{C}{\tau_j},
  \label{eq:epoch-volume}
\end{equation}
and that $H=C_0/\sqrt\alpha$ restricted iterations reduce the relevant error
by a constant factor.  If discovery and transition between epochs introduce
no additional full scans, then
\begin{align}
 T&=\bigO\!\left(\frac1{\sqrt\alpha}
                  \log\frac{\tau_0}{\tau_J}\right),
 \label{eq:epoch-iterations}\\
 \mathcal W
 &\leq H\sum_{j=0}^{J}\vol(U_j)
 =\bigO\!\left(\frac1{\sqrt\alpha\,\tau_J}\right).
 \label{eq:epoch-work}
\end{align}
\end{proposition}

\begin{proof}
The iteration bound is $H(J+1)$.  Because $1/\tau_j$ grows geometrically,
$\sum_j1/\tau_j=\bigO(1/\tau_J)$, which gives \cref{eq:epoch-work}.
\end{proof}

The proposition is only an accounting implication.  It assumes precisely the
two properties that need proof: a locally discoverable certified region and
stable acceleration as that region grows.

\subsection{Why constant-frequency restarting loses acceleration}

For a degree-$H$ polynomial method on an SPD spectrum $[\alpha,1]$, the best
worst-case block contraction is
\begin{equation}
  \rho_H^\star
  =\frac1{\abs{T_H((1+\alpha)/(1-\alpha))}},
  \label{eq:cheb-block}
\end{equation}
where $T_H$ is the Chebyshev polynomial of the first kind.  If
$H\sqrt\alpha\ll1$, then
\begin{equation}
  1-\rho_H^\star=\Theta(H^2\alpha).
  \label{eq:short-block}
\end{equation}
Repeating such blocks therefore needs
\begin{equation}
  \Omega\!\left(\frac1{H\alpha}\log\frac1\epsppr\right)
  \label{eq:restart-lower}
\end{equation}
total polynomial degree.  In particular, $H=\bigO(1)$ returns to the
nonaccelerated $\alpha^{-1}$ scale.  An accelerated accuracy epoch must retain
momentum for $H=\Omega(1/\sqrt\alpha)$ steps.  This classical polynomial fact
does not by itself resolve support changes, but it rules out constant-step
ordinary restarts as the flattening mechanism.

\section{A canonical exact PPR core}
\label{sec:exact-core}

Before constructing an implementable gate, it is useful to see that a small
accurate region always exists.  Define normalized PPR values
\begin{equation}
  y^0:=D^{-1/2}x^0=D^{-1}\pi
\end{equation}
and the exact superlevel set
\begin{equation}
  U_\tau^0:=\{i:y_i^0>\tau\}.
  \label{eq:exact-core}
\end{equation}

\begin{theorem}[Exact superlevel core]
\label{thm:exact-core}
The core in \cref{eq:exact-core} satisfies
\begin{equation}
  \vol(U_\tau^0)\leq\frac1\tau.
  \label{eq:exact-core-volume}
\end{equation}
Let $x^{U}$ solve the Dirichlet restriction
\begin{equation}
  Q_{UU}x_U^U=b_U,
  \qquad x_{U^c}^U=0,
  \qquad U=U_\tau^0.
  \label{eq:dirichlet}
\end{equation}
Then
\begin{equation}
  0\leq x^U\leq x^0,
  \qquad
  \norm{D^{-1/2}(x^0-x^U)}_\infty\leq\tau.
  \label{eq:dirichlet-error}
\end{equation}
\end{theorem}

\begin{proof}
By \cref{eq:ppr-mass},
$\sum_i d_i y_i^0=1$.  Markov's inequality gives
\cref{eq:exact-core-volume}.  Positivity of $Q_{UU}^{-1}$ and the fact that
$Q_{UU}x_U^0=b_U-Q_{UU^c}x_{U^c}^0\geq b_U$ give
$0\leq x^U\leq x^0$.

Set $e=x^0-x^U$ and $h=D^{-1/2}e$.  For $i\in U$,
$(Qe)_i=0$, hence
\begin{equation}
  h_i=\frac{1-\alpha}{1+\alpha}
      \sum_j\frac{A_{ij}}{d_i}h_j.
  \label{eq:discounted-harmonic}
\end{equation}
On $U^c$, $0\leq h_i=y_i^0\leq\tau$.  If a positive maximum larger than
$\tau$ occurred in $U$, \cref{eq:discounted-harmonic} would make it strictly
smaller than a convex combination of values no larger than that maximum, a
contradiction.  This proves \cref{eq:dirichlet-error}.
\end{proof}

The theorem is existential because $U_\tau^0$ depends on the unknown exact
solution.  It nevertheless shows that a peak volume larger than
$\Theta(1/\epsppr)$ is not information-theoretically required by the output
criterion \cref{eq:note-error}.

\subsection{Why a residual wave can be unnecessary on a spider}

For the center-source spider $G_{B,N}$, direct solution of the radial
recurrence gives
\begin{equation}
  y_o^0
  =\frac{\sqrt\alpha}{B}
    \frac{1+\chi^{2N}}{1-\chi^{2N}},
  \qquad
  \chi:=\frac{1-\sqrt\alpha}{1+\sqrt\alpha}.
  \label{eq:spider-center-value}
\end{equation}
For long arms, $y_o^0\to\sqrt\alpha/B$.  If
$B=1/(4\epsppr)$ and $\alpha<1/16$, then the limiting center value is below
$\epsppr$, and all arm values are smaller.  The true $\epsppr$-important
core is empty even though a residual-based process may launch a wave from the
center.  This example motivates a gate based on solution/KKT importance rather
than momentum residual magnitude.

\section{RPPR supplies an implementable safe core}
\label{sec:rppr-core}

Let $S_\rho:=\supp(x_\rho^\star)$.  The following properties turn fixed
regularization into a localization certificate.

\begin{theorem}[RPPR support and approximation]
\label{thm:rppr-support}
For every $\rho>0$,
\begin{align}
  \vol(S_\rho)&\leq\frac1\rho,
  \label{eq:rppr-volume}\\
  0\leq D^{-1/2}(x^0-x_\rho^\star)&\leq\rho\one.
  \label{eq:rppr-ppr-error}
\end{align}
Moreover, if $\rho'\geq\rho$, then
$x_{\rho'}^\star\leq x_\rho^\star$ and
$S_{\rho'}\subseteq S_\rho$.
\end{theorem}

\begin{proof}
Regularization-path monotonicity is the source result stated, for example, in
the appendix of \citet{fountoulakis2026}.  We prove the two quantitative
claims.  Let $\zeta_\rho:=b-Qx_\rho^\star$.  On $S_\rho$, KKT gives
$\zeta_{\rho,i}=\alpha\rho\sqrt{d_i}$.  Outside $S_\rho$,
$x_{\rho,i}^\star=0$ and the off-diagonal entries of $Q$ are nonpositive, so
$(Qx_\rho^\star)_i\leq0$ and hence $\zeta_{\rho,i}\geq0$.  The full
$\ell_1$ KKT condition also gives
$\zeta_\rho\leq\alpha\rho D^{1/2}\one$.  Therefore
\begin{align*}
 \alpha\rho\vol(S_\rho)
 &\leq\one^\top D^{1/2}\zeta_\rho\\
 &=\one^\top D^{1/2}b
   -\one^\top D^{1/2}Qx_\rho^\star\\
 &=\alpha-\alpha\one^\top D^{1/2}x_\rho^\star
 \leq\alpha,
\end{align*}
which proves \cref{eq:rppr-volume}.  Finally,
\[
  x^0-x_\rho^\star=Q^{-1}\zeta_\rho,
\]
and \cref{eq:Q-properties} gives
\[
  0\leq D^{-1/2}Q^{-1}\zeta_\rho
  \leq\alpha\rho D^{-1/2}Q^{-1}D^{1/2}\one
  =\rho\one.
\]
\end{proof}

Thus RPPR is not merely a different target: when $\rho$ is chosen on the
desired error scale, it is an internal localization device for unregularized
PPR.

\section{Safe envelopes for signed accelerated candidates}
\label{sec:safe-envelope}

Acceleration may overshoot, so its candidate cannot itself control support.
The next lemma converts any signed restricted candidate into a certified lower
point.

For $U\subseteq S_\rho$, let $x_U^\star$ solve the restricted nonnegative
quadratic program
\begin{equation}
  x_U^\star\in\arg\min_{z_U\geq0}
  \left\{\frac12z_U^\top Q_{UU}z_U-c_{\rho,U}^\top z_U\right\},
  \qquad x_{U^c}^\star=0.
  \label{eq:restricted-QP}
\end{equation}
The Stieltjes obstacle comparison principle gives
\begin{equation}
  0\leq x_U^\star\leq x_\rho^\star.
  \label{eq:restricted-comparison}
\end{equation}
Indeed, $x_{\rho,U}^\star$ is a supersolution of the restricted KKT system
because $U\subseteq S_\rho$ and
$Q_{UU}x_{\rho,U}^\star-c_{\rho,U}
=-Q_{UU^c}x_{\rho,U^c}^\star\geq0$.

\begin{theorem}[Signed candidate to safe lower envelope]
\label{thm:safe-envelope}
Let $z_U\in\R^U$ be arbitrary and define
\begin{equation}
 q_U:=c_{\rho,U}-Q_{UU}z_U,
 \qquad
 \delta_U:=\frac1\alpha
  \norm{D_U^{-1/2}(q_U)_-}_\infty.
 \label{eq:envelope-delta}
\end{equation}
Set
\begin{equation}
 \ell_U:=\left[z_U-\delta_UD_U^{1/2}\one\right]_+,
 \qquad \ell_{U^c}:=0.
 \label{eq:safe-envelope}
\end{equation}
Then
\begin{equation}
  0\leq\ell\leq x_U^\star\leq x_\rho^\star.
  \label{eq:envelope-order}
\end{equation}
\end{theorem}

\begin{proof}
Let $q_U^\star:=c_{\rho,U}-Q_{UU}x_U^\star\leq0$ be the restricted KKT
residual.  Since $Q_{UU}^{-1}\geq0$,
\begin{align*}
 x_U^\star-z_U
 &=Q_{UU}^{-1}(q_U-q_U^\star)
 \geq-Q_{UU}^{-1}(q_U)_-.
\end{align*}
The full eigenvector identity in \cref{eq:Q-properties} and the signs of the
cut entries imply
\begin{equation}
 Q_{UU}D_U^{1/2}\one
 =\alpha D_U^{1/2}\one-Q_{UU^c}D_{U^c}^{1/2}\one
 \geq\alpha D_U^{1/2}\one.
\end{equation}
Thus
$Q_{UU}^{-1}D_U^{1/2}\one\leq\alpha^{-1}D_U^{1/2}\one$.
By \cref{eq:envelope-delta},
$Q_{UU}^{-1}(q_U)_-\leq\delta_UD_U^{1/2}\one$, and hence
$z_U-\delta_UD_U^{1/2}\one\leq x_U^\star$.  Taking the positive part preserves
the inequality because $x_U^\star\geq0$.  The last comparison is
\cref{eq:restricted-comparison}.
\end{proof}

The accelerator is therefore free to use signed momentum in $z$; only
$\ell$ is allowed to govern expansion.

\subsection{Safe admission and a one-sided stopping certificate}

\begin{lemma}[Safe admission]
\label{lem:safe-admission}
Let $\ell\leq x_\rho^\star$ be nonnegative.  If $\ell_i=0$ and
\begin{equation}
  \kappa_{\rho,i}(\ell)<0,
  \label{eq:negative-violation}
\end{equation}
then $i\in S_\rho$.
\end{lemma}

\begin{proof}
If $i\notin S_\rho$, then $\ell_i=x_{\rho,i}^\star=0$.  Since all
off-diagonal entries in row $i$ of $Q$ are nonpositive and
$\ell\leq x_\rho^\star$,
\[
 \kappa_{\rho,i}(\ell)-\kappa_{\rho,i}(x_\rho^\star)
 =\sum_{j\neq i}Q_{ij}(\ell_j-x_{\rho,j}^\star)\geq0.
\]
KKT gives $\kappa_{\rho,i}(x_\rho^\star)\geq0$, contradicting
\cref{eq:negative-violation}.
\end{proof}

\begin{lemma}[One-sided maximum-principle certificate]
\label{lem:one-sided-certificate}
Let $0\leq\ell\leq x_\rho^\star$.  If
\begin{equation}
  \kappa_{\rho,i}(\ell)\geq-\alpha\tau\sqrt{d_i}
  \qquad\text{for every }i,
  \label{eq:global-certificate}
\end{equation}
then
\begin{equation}
  \norm{D^{-1/2}(x_\rho^\star-\ell)}_\infty\leq\tau.
  \label{eq:certificate-error}
\end{equation}
\end{lemma}

\begin{proof}
Let $e=x_\rho^\star-\ell\geq0$, and choose $i$ maximizing
$M:=e_i/\sqrt{d_i}$.  If $M=0$ there is nothing to prove.  Otherwise
$x_{\rho,i}^\star>0$, so $\kappa_{\rho,i}(x_\rho^\star)=0$.  The normalized
maximum principle gives
\begin{align*}
 \frac{(Qe)_i}{\sqrt{d_i}}
 &=\frac{1+\alpha}{2}M
   -\frac{1-\alpha}{2}
     \sum_j\frac{A_{ij}}{d_i}\frac{e_j}{\sqrt{d_j}}
 \geq\alpha M.
\end{align*}
On the other hand,
$(Qe)_i=\kappa_{\rho,i}(x_\rho^\star)-\kappa_{\rho,i}(\ell)
=-\kappa_{\rho,i}(\ell)$.  Condition \cref{eq:global-certificate} yields
$\alpha M\leq\alpha\tau$.
\end{proof}

\begin{corollary}[PPR accuracy and peak volume]
\label{cor:accuracy-volume}
Suppose a process maintains $U_t\subseteq S_\rho$, constructs $\ell_t$ by
\cref{eq:safe-envelope}, admits only vertices satisfying
\begin{equation}
  \kappa_{\rho,i}(\ell_t)<-\alpha\tau\sqrt{d_i},
  \label{eq:strong-admission}
\end{equation}
and stops when \cref{eq:global-certificate} holds.  Then
\begin{equation}
  \max_t\vol(U_t)\leq\frac1\rho,
  \qquad
  \norm{D^{-1/2}(x^0-\ell_t)}_\infty\leq\rho+\tau.
  \label{eq:gate-main-bound}
\end{equation}
In particular, $\rho=\tau=\epsppr/2$ gives peak volume at most
$2/\epsppr$ and PPR error at most $\epsppr$.
\end{corollary}

\begin{proof}
Safe admission follows from \cref{lem:safe-admission}, and the volume bound
then follows from \cref{eq:rppr-volume}.  The error bound is the triangle
inequality applied to \cref{eq:rppr-ppr-error} and
\cref{eq:certificate-error}.
\end{proof}

Only $U_t$, its graph boundary, and source coordinates need inspection.  If a
vertex is outside $U_t\cup N(U_t)$, then $(Q\ell_t)_i=0$; if its initial
gradient was strongly negative it would already have been admitted, and
otherwise it automatically satisfies the tolerance in
\cref{eq:global-certificate}.  A boundary check is therefore chargeable to
the incident edges of the current certified set, subject to the usual sparse
source assumption.

\subsection{An exact stagewise gate is correct but not yet fast}

Consider the idealized procedure that begins with $U=\varnothing$, exactly
solves \cref{eq:restricted-QP}, adds every outside vertex satisfying
\cref{eq:strong-admission}, and repeats.  At an exact interior restricted
solution $x_U^\star$, the envelope correction is zero.  If $B$ is the next
batch, $h_B:=\kappa_{\rho,B}(x_U^\star)<0$.  The Schur-complement formula used in
\cref{sec:shock} below shows that the enlarged unconstrained minimizer has
positive correction
\begin{equation}
 w_B=-S_B^{-1}h_B>0,
 \qquad
 w_U=-Q_{UU}^{-1}Q_{UB}w_B\geq0.
 \label{eq:positive-expansion}
\end{equation}
Thus every admitted coordinate is positive at the enlarged restricted
minimizer, and induction justifies the interior claim.  Every batch lies in
$S_\rho$ by safe admission.  If no batch exists, the exact restricted
gradient vanishes on $U$ and the global certificate holds outside $U$.
Consequently the stagewise exact process terminates after at most
$|S_\rho|$ nonempty admissions and satisfies \cref{eq:gate-main-bound}.

This proves correctness and peak-volume flattening.  Exactly re-solving after
every batch can still repeat too much work; acceleration across expansions is
the remaining issue.

\section{A path refutes the logarithmic-restart claim}
\label{sec:path}

The initial proposal hoped to group all support changes into only
$\bigO(\log(1/\epsppr))$ ordinary restarts.  A path rules this out.

Let $P_N=(v_0,v_1,\ldots,v_N)$ be a path with source $s=e_{v_0}$.  Write
$y_i^0=x_i^0/\sqrt{d_i}$.  Solving the recurrence gives
\begin{equation}
 y_i^0
 =\sqrt\alpha\,
   \frac{\chi^i+\chi^{2N-i}}{1-\chi^{2N}}
 \geq\sqrt\alpha\,\chi^i,
 \qquad
 \chi=\frac{1-\sqrt\alpha}{1+\sqrt\alpha}.
 \label{eq:path-solution}
\end{equation}
By \cref{eq:rppr-ppr-error},
\begin{equation}
  \frac{x_{\rho,i}^\star}{\sqrt{d_i}}
  \geq\sqrt\alpha\,\chi^i-\rho.
  \label{eq:path-rppr-lower}
\end{equation}
Define
\begin{equation}
 J_{\mathrm{path}}(\rho):=\left\lfloor
   \frac{\log(\sqrt\alpha/(3\rho))}{-\log\chi}
 \right\rfloor,
 \label{eq:path-radius}
\end{equation}
assuming $3\rho<\sqrt\alpha$ and $N>J_{\mathrm{path}}(\rho)$.

\begin{theorem}[One-coordinate expansion on a path]
\label{thm:path-expansion}
Run the exact safe gate with $\tau=\rho$.  Once
$U_j=\{v_0,\ldots,v_j\}$, every stage with
$j<J_{\mathrm{path}}(\rho)$ admits precisely the
next vertex $v_{j+1}$ and no more distant vertex.  Consequently any method
that restarts after every safe expansion makes at least
\begin{equation}
 \Omega\!\left(
   \frac1{\sqrt\alpha}
   \log\frac{\sqrt\alpha}{\rho}
 \right)
 \label{eq:path-restarts}
\end{equation}
restarts for $0<\alpha\leq1/4$ and sufficiently small $\rho$.
\end{theorem}

\begin{proof}
Let $x^{U_j}$ be the exact restricted minimizer.  For $i\geq j+2$, vertex
$v_i$ has no neighbor in $U_j$ and is not the source, so
\begin{equation}
  \kappa_{\rho,i}(x^{U_j})=\alpha\rho\sqrt{d_i}>0.
  \label{eq:path-far-gradient}
\end{equation}
On the other hand, \cref{eq:path-rppr-lower} gives normalized RPPR value at
least $2\rho$ through radius $J_{\mathrm{path}}(\rho)$. Since
$x^{U_j}_{j+1}=0$, the
restricted point is more than $\rho$ away from $x_\rho^\star$ in normalized
infinity norm.  The contrapositive of \cref{lem:one-sided-certificate} implies
that some coordinate has gradient below $-\alpha\rho\sqrt{d_i}$.  The exact
restricted gradient vanishes on $U_j$, while
\cref{eq:path-far-gradient} rules out every $i\geq j+2$.  Hence only
$v_{j+1}$ can violate, proving the first claim.  Finally,
\begin{equation}
  -\log\chi
  =2\operatorname{artanh}(\sqrt\alpha)
  =\Theta(\sqrt\alpha)
\end{equation}
for $\alpha\leq1/4$, which yields \cref{eq:path-restarts}.
\end{proof}

This lower bound is exactly on the accelerated propagation scale.  It does
not refute restart-free continuation; it refutes only the assertion that the
number of support-triggered ordinary restarts is independent of the
$1/\sqrt\alpha$ propagation distance.

\section{Expansion shock: exact identity and corrected direction}
\label{sec:shock}

Let $x_U^\star$ be an interior exact restricted minimizer and let $B$ be a
batch of negative-gradient coordinates.  Define
\begin{equation}
 h_B:=\kappa_{\rho,B}(x_U^\star),
 \qquad
 S_B:=Q_{BB}-Q_{BU}Q_{UU}^{-1}Q_{UB}.
 \label{eq:schur}
\end{equation}
Assume the enlarged optimum is interior, as it is for the negative-gradient
admissions in the exact stagewise process.

\begin{proposition}[Exact expansion gain]
\label{prop:shock}
The improvement in the restricted optimal value is
\begin{equation}
 \Delta_B
 :=G_\rho(x_U^\star)-G_\rho(x_{U\cup B}^\star)
 =\frac12h_B^\top S_B^{-1}h_B.
 \label{eq:shock-exact}
\end{equation}
Moreover,
\begin{equation}
  \alpha I\preceq S_B\preceq I,
  \qquad
  \frac12\norm{h_B}_2^2
  \leq\Delta_B
  \leq\frac1{2\alpha}\norm{h_B}_2^2.
  \label{eq:shock-bounds}
\end{equation}
\end{proposition}

\begin{proof}
Eliminating the old variables from the enlarged quadratic gives the reduced
problem
$\min_w\{\frac12w^\top S_Bw+h_B^\top w\}$, whose minimizer is
$w=-S_B^{-1}h_B$ and whose gain is \cref{eq:shock-exact}.  The Schur
complement is the minimum of the full quadratic over the eliminated block;
using $\alpha I\preceq Q\preceq I$ gives
$\alpha I\preceq S_B\preceq I$.  Inverting these inequalities yields
\cref{eq:shock-bounds}.
\end{proof}

The lower bound in \cref{eq:shock-bounds} is useful for proving that an
expansion makes genuine progress.  It cannot upper-bound the perturbation of
an accelerated Lyapunov function.  The upper bound may be larger by
$1/\alpha$.  This corrects the earlier proposal to restart only when a sum of
lower shock bounds reaches a fixed energy fraction.

\section{What telescopes: an amortized re-solving lemma}
\label{sec:amortized}

Although the number of expansions can be large, the continuous optimization
work needed to adapt to them admits a clean amortization.

For a safe point $0\leq x\leq x_\rho^\star$, define the one-sided violation
\begin{equation}
  \nu(x):=\norm{D^{-1/2}(\kappa_\rho(x))_-}_\infty.
  \label{eq:nu}
\end{equation}

\begin{lemma}[Violation to objective gap]
\label{lem:violation-gap}
For every safe point,
\begin{align}
 \norm{D^{-1/2}(x_\rho^\star-x)}_\infty
 &\leq\frac{\nu(x)}\alpha,
 \label{eq:nu-error}\\
 \nu(x)\leq\alpha\sigma
 \quad&\Longrightarrow\quad
 G_\rho(x)-G_\rho(x_\rho^\star)
 \leq\frac{\sigma^2}{2\rho}.
 \label{eq:nu-gap}
\end{align}
\end{lemma}

\begin{proof}
The proof of \cref{lem:one-sided-certificate} with
$\tau=\nu(x)/\alpha$ gives \cref{eq:nu-error}.  Put
$e=x_\rho^\star-x$.  Because $e$ is supported on $S_\rho$ and complementarity
annihilates the linear term,
\begin{equation}
 G_\rho(x)-G_\rho(x_\rho^\star)=\frac12e^\top Qe
 \leq\frac12\norm{e}_2^2.
\end{equation}
If $\nu(x)\leq\alpha\sigma$, then $e_i\leq\sigma\sqrt{d_i}$ by
\cref{eq:nu-error}.  Hence
\[
 \norm{e}_2^2
 \leq\sigma^2\vol(S_\rho)
 \leq\frac{\sigma^2}\rho,
\]
using \cref{eq:rppr-volume}.
\end{proof}

\begin{assumption}[Restricted accelerated contraction]
\label{ass:contraction}
On every fixed restricted QP, $k$ full restricted iterations reduce objective
error by a factor at most $\exp(-k/H)$, where
$H=C/\sqrt\alpha$ and $C$ is universal.
\end{assumption}

Strongly convex accelerated proximal gradient has this fixed-subspace
property.  The assumption says nothing about continuing momentum across a
subspace change.

\begin{theorem}[Amortized re-solving]
\label{thm:amortized}
Let $\sigma_{j-1}=2\sigma_j$ be geometric accuracy epochs and put
\begin{equation}
  E_j:=\frac{\sigma_j^2}{2\rho}.
\end{equation}
Suppose epoch $j$ begins at a safe point with
$\nu(x)\leq\alpha\sigma_{j-1}$, and restore restricted objective error to
$E_j$ after every support expansion.  Under \cref{ass:contraction}, the total
number of restricted iterations over all epochs down to $\sigma_J=\tau$ is
\begin{equation}
 K
 \leq\bigO\!\left(
       \frac1{\sqrt\alpha}\log\frac{\sigma_0}{\tau}
      \right)+N_{\mathrm{exp}},
 \label{eq:amortized-iterations}
\end{equation}
where $N_{\mathrm{exp}}$ is the number of nonempty expansion events.  Since
$\vol(U)\leq1/\rho$, full-sweep work obeys
\begin{equation}
 \mathcal W
 \leq\bigO\!\left(
       \frac1{\rho\sqrt\alpha}\log\frac{\sigma_0}{\tau}
      \right)
      +\bigO\!\left(\frac{N_{\mathrm{exp}}}{\rho}\right).
 \label{eq:amortized-work}
\end{equation}
\end{theorem}

\begin{proof}
By \cref{lem:violation-gap}, the global objective gap at the start of epoch
$j$ is at most $4E_j$.  Restricted optimal values decrease monotonically
toward the global value, so all expansion gains $\Delta_{j,\ell}$ in that
epoch telescope:
\begin{equation}
  \sum_\ell\Delta_{j,\ell}\leq4E_j.
  \label{eq:shock-telescope}
\end{equation}
If the pre-expansion restricted error is at most $E_j$, then immediately after
an optimum shift of size $\Delta_{j,\ell}$ it is at most
$E_j+\Delta_{j,\ell}$.  By \cref{ass:contraction}, restoring it to $E_j$
takes at most
\begin{equation}
 k_{j,\ell}
 =\left\lceil H\log\left(1+\frac{\Delta_{j,\ell}}{E_j}\right)\right\rceil.
\end{equation}
Therefore
\begin{align*}
 \sum_\ell k_{j,\ell}
 &\leq H\sum_\ell
       \log\left(1+\frac{\Delta_{j,\ell}}{E_j}\right)
       +N_{\mathrm{exp},j}\\
 &\leq\frac{H}{E_j}\sum_\ell\Delta_{j,\ell}
       +N_{\mathrm{exp},j}
 \leq4H+N_{\mathrm{exp},j}.
\end{align*}
Summing over $\bigO(\log(\sigma_0/\tau))$ epochs proves
\cref{eq:amortized-iterations}.  Multiplying full iterations by the uniform
volume cap proves \cref{eq:amortized-work}.
\end{proof}

This theorem improves over independently re-solving every enlarged subspace:
all continuous shock sizes telescope.  The unresolved term is discrete.  In
the worst case $N_{\mathrm{exp}}$ may be as large as the support size, and the
proof above charges a full current-support scan to every such event.

\section{The correct open lemma: one-sided expanding-subspace acceleration}
\label{sec:open-lemma}

The path in \cref{sec:path} shows that support may expand on nearly every
accelerated-scale propagation step.  The desired algorithm must therefore
continue through expansions, or charge each expansion only to newly exposed
edges.

A concrete candidate zero-pads its history when $U_t$ expands and otherwise
uses the standard strongly convex momentum
\begin{align}
 y^{(t)}
 &=x^{(t)}+\beta_{\mathrm F}\bigl(x^{(t)}-x^{(t-1)}\bigr),
 \qquad
 \beta_{\mathrm F}:=\frac{1-\sqrt\alpha}{1+\sqrt\alpha},
 \label{eq:dynamic-y}\\
 z^{(t+1)}
 &=\Pi_{\R_+^{U_t}}
   \bigl(y^{(t)}-(Qy^{(t)}-c_\rho)\bigr),
 \label{eq:dynamic-step}
\end{align}
where vectors are zero outside $U_t$.  Construct $\ell^{(t+1)}$ from
$z^{(t+1)}$ by \cref{eq:envelope-delta,eq:safe-envelope}, then set
\begin{equation}
 U_{t+1}:=U_t\cup
 \left\{i:
 \kappa_{\rho,i}(\ell^{(t+1)})
 <-\alpha\tau\sqrt{d_i}
 \right\}.
 \label{eq:dynamic-gate}
\end{equation}
The proved gate invariants hold regardless of whether the following
conjecture is true:
\begin{equation}
 \ell^{(t)}\leq x_\rho^\star,
 \qquad
 U_t\subseteq S_\rho,
 \qquad
 \vol(U_t)\leq\frac1\rho.
 \label{eq:dynamic-invariants}
\end{equation}

\begin{conjecture}[One-sided expanding-subspace acceleration]
\label{conj:one-sided-continuation}
For the recurrence \cref{eq:dynamic-y,eq:dynamic-step,eq:dynamic-gate}, or for
a closely related projected Chebyshev recurrence, universal constants $C,c>0$
exist such that
\begin{equation}
 \norm{D^{-1/2}(x_\rho^\star-\ell^{(t)})}_\infty
 \leq C\exp(-c\sqrt\alpha\,t)+\tau.
 \label{eq:desired-continuation}
\end{equation}
\end{conjecture}

If \cref{conj:one-sided-continuation} holds, then
\begin{equation}
 t=\bigO\!\left(\frac1{\sqrt\alpha}\log\frac1\tau\right),
 \qquad
 \mathcal W
 =\bigO\!\left(
   \frac1{\rho\sqrt\alpha}\log\frac1\tau
 \right).
 \label{eq:conditional-final}
\end{equation}
With $\rho=\tau=\epsppr/2$, this is the sought
$\softO(1/(\sqrt\alpha\epsppr))$ bound.

\paragraph{Why existing analyses do not immediately prove it.}
The nonhomogeneous \LOCCH{} formula of \citet{zhou2024} controls a full signed
inactive residual.  The safe gate controls only the negative part of the KKT
residual at a corrected lower envelope.  Outside the optimal support, positive
KKT slack can be dense and is harmless for the obstacle problem, but it enters
a two-sided residual norm.  Energy contraction alone also does not imply the
uniform active-volume invariant.  Finally, \cref{prop:shock} supplies an
upper perturbation bound only with a possible $1/\alpha$ loss.

\paragraph{A weaker sufficient target.}
It is enough to strengthen \cref{thm:amortized} so that discrete expansion
overhead is charged to newly admitted volume:
\begin{equation}
  \mathcal W_{\mathrm{exp}}
  \leq C\sum_t\vol(U_{t+1}\setminus U_t)
  \leq\frac{C}{\rho},
  \label{eq:new-volume-charge}
\end{equation}
rather than $N_{\mathrm{exp}}/\rho$.  Combining
\cref{eq:new-volume-charge} with the continuous telescoping term already gives
the target up to a logarithm.  This newly-admitted-edge lemma may be easier
than proving the pointwise contraction \cref{eq:desired-continuation}.

\section{Implications for the hybrid-solver program}
\label{sec:implications}

The conversation leads to four firm conclusions.

\begin{enumerate}[leftmargin=2em]
\item \emph{Acceleration has not vanished spectrally.}
The spider's restricted condition number is never worse than $1/\alpha$.
The extra factor observed in local work comes from the space--time footprint
of repeated scans.

\item \emph{Peak-volume flattening is rigorously feasible.}
RPPR KKT structure supplies a locally checkable, momentum-independent gate
whose working set never exceeds volume $1/\rho$.  This part no longer depends
on a conjectural signed-residual $\ell_1$ invariant.

\item \emph{Naive restart and batching mechanisms are insufficient.}
Short fixed restart blocks lose acceleration; FIFO batch splitting preserves
triangular work; and a path forces accelerated-scale many safe expansions.

\item \emph{The open problem is now narrow.}
One must prove one-sided accelerated tracking over nested obstacle subspaces,
or amortize expansion work by newly exposed volume.  Until then, the
graph-uniform $\softO(1/(\sqrt\alpha\epsppr))$ theorem remains open.
\end{enumerate}

The gate is compatible with the broader hybrid philosophy: use signed AESP,
Chebyshev, SOR, or another accelerator to generate candidates, but let a
monotone KKT envelope decide which coordinates may be exposed.  A
momentum-free local cleanup remains a proof-safe fallback.  The new result
does not close the separate early-AESP locality promotion gate in the
repository's AESP--\LOCSOR{} note; it instead supplies an alternative route
whose own continuation lemma must be resolved.

\subsection{Recommended proof and falsification sequence}

\begin{enumerate}[leftmargin=2em]
\item Prove \cref{eq:desired-continuation} first on endpoint-source paths,
where support expands by one coordinate and the recurrence is scalar after a
Schur complement.
\item Extend the proof to trees with bounded branching, isolating whether
simultaneous boundary admissions create an unavoidable $1/\alpha$ shock.
\item Test the zero-padded recurrence on adversarial spiders, stars of stars,
and lollipops.  Record peak volume, newly admitted volume, edge revisits,
$\nu(\ell_t)$, and objective shock separately.
\item If pointwise continuation fails, target the weaker charge
\cref{eq:new-volume-charge}; if that also fails, construct the corresponding
M-matrix counterexample and state the structural condition that restores it.
\end{enumerate}

\section{Claim ledger}
\label{sec:ledger}

\begin{center}
\begin{tabularx}{\textwidth}{@{}X p{0.22\textwidth}@{}}
\toprule
Claim & Status \\
\midrule
$\kappa(Q_{UU})\leq1/\alpha$ for every working set & Proved,
\cref{lem:principal-conditioning} \\
Exact depth-$L$ spider condition number is
$\Theta(1/(\alpha+L^{-2}))$ & Proved, \cref{prop:spider-spectrum} \\
An $\epsppr$-accurate PPR Dirichlet core of volume $1/\epsppr$ exists &
Proved, \cref{thm:exact-core} \\
Fixed RPPR regularization supplies a safe $1/\rho$ support & Proved/source,
\cref{thm:rppr-support} \\
Signed accelerated candidates can be used without unsafe admissions & Proved,
\cref{thm:safe-envelope,lem:safe-admission} \\
The exact safe gate has peak volume $2/\epsppr$ and PPR error $\epsppr$ &
Proved, \cref{cor:accuracy-volume} \\
Only $\bigO(\log(1/\epsppr))$ support restarts are needed & Refuted by path,
\cref{thm:path-expansion} \\
Expansion shocks can be absorbed using only their elementary lower bound &
Refuted; direction corrected in \cref{prop:shock} \\
Continuous restricted re-solving has accelerated amortized cost & Proved
under fixed-subspace contraction, \cref{thm:amortized} \\
The dynamic gate achieves
$\softO(1/(\rho\sqrt\alpha))$ graph-uniform work & Open,
\cref{conj:one-sided-continuation} \\
\bottomrule
\end{tabularx}
\end{center}

\begin{thebibliography}{99}

\bibitem[Andersen et al.(2007)Andersen, Chung, and Lang]{andersen2007}
Reid Andersen, Fan Chung, and Kevin Lang.
\newblock Using PageRank to locally partition a graph.
\newblock \emph{Internet Mathematics}, 4(1):35--64, 2007.

\bibitem[Fountoulakis et al.(2019)Fountoulakis, Roosta-Khorasani, Shun,
Cheng, and Mahoney]{fountoulakis2019}
Kimon Fountoulakis, Farbod Roosta-Khorasani, Julian Shun, Xiang Cheng, and
Michael Mahoney.
\newblock Variational perspective on local graph clustering.
\newblock \emph{Mathematical Programming}, 174:553--573, 2019.

\bibitem[Fountoulakis and Martínez-Rubio(2026)]{fountoulakis2026}
Kimon Fountoulakis and David Martínez-Rubio.
\newblock Complexity of classical acceleration for
$\ell_1$-regularized PageRank.
\newblock arXiv:2602.21138v2, 2026.

\bibitem[Huang et al.(2025)Huang, Luo, Xiao, Yang, and Zhou]{huang2025}
Binbin Huang, Luo Luo, Yanghua Xiao, Deqing Yang, and Baojian Zhou.
\newblock Accelerated evolving set processes for local PageRank computation.
\newblock \emph{Advances in Neural Information Processing Systems}, 38, 2025.

\bibitem[Lin et al.(2018)Lin, Mairal, and Harchaoui]{lin2018}
Hongzhou Lin, Julien Mairal, and Zaid Harchaoui.
\newblock Catalyst acceleration for first-order convex optimization: From
theory to practice.
\newblock \emph{Journal of Machine Learning Research}, 18(212):1--54, 2018.

\bibitem[Martínez-Rubio et al.(2023)Martínez-Rubio, Wirth, and
Pokutta]{martinezrubio2023}
David Martínez-Rubio, Elias Wirth, and Sebastian Pokutta.
\newblock Accelerated and sparse algorithms for approximate personalized
PageRank and beyond.
\newblock In \emph{Proceedings of COLT}, volume 195 of PMLR, pages
2852--2876, 2023.

\bibitem[O'Donoghue and Candès(2015)]{odonoghue2015}
Brendan O'Donoghue and Emmanuel Candès.
\newblock Adaptive restart for accelerated gradient schemes.
\newblock \emph{Foundations of Computational Mathematics}, 15:715--732, 2015.

\bibitem[Zhou et al.(2024)Zhou, Sun, Babanezhad Harikandeh, Guo, Yang, and
Xiao]{zhou2024}
Baojian Zhou, Yifan Sun, Reza Babanezhad Harikandeh, Xingzhi Guo, Deqing Yang,
and Yanghua Xiao.
\newblock Iterative methods via locally evolving set process.
\newblock \emph{Advances in Neural Information Processing Systems}, 37, 2024.

\end{thebibliography}

\end{document}
